\begin{frame}{스펙트럼 읽는 법 --- 시험 답안 구조}
  \begin{itemize}
    \item \kb{왼쪽 $\to$ 오른쪽으로 갈수록 모델의 유연성(가용
      자유도)이 커진다} --- Ridge는 10개 가중치에 $\lambda$로
      매달아 자유도를 조절하는 ``연속 스펙트럼''의 한 점,
      무제한 트리는 265건을 전부 외우는 ``스펙트럼의 끝''.
    \item 격차가 큰 모델(오른쪽)은 \textbf{분산}이 크다
      (Week 4.1: ``모델 추정량이 train 데이터의 선택에 따라
      크게 흔들린다''). 격차가 작은 모델(왼쪽)은 \textbf{편향}
      이 크다 (Week 6: ``평균적으로 진실을 향한 systematic한
      오차'').
  \end{itemize}
  \vskip0.4em
  \begin{block}{트리의 val R$^2$가 \textbf{0 아래}(음수)}
    train에서 1.000을 찍고도 검증 데이터에서는 \textbf{평균만
    예측하는 모델보다 나쁘다}는 뜻 --- Week 6.2 U자 곡선의
    ``$\lambda\to0$ 쪽 끝''이 실제로 이렇게 처져 있음을
    보여주는 숫자. \kq{``train R$^2$=1.000이니까 좋은 모델''
    이라는 답안은 이 한 줄로 반박된다.}
  \end{block}
  \vskip0.3em
  {\footnotesize Ridge($10^{-4}$)와 Ridge(1)가 거의 같은 이유:
  표준화한 265개 샘플에서 $\lambda=1$은 데이터 스케일 대비
  \textbf{아직 매우 작은 벌금} --- ``$\lambda$가 1로 작아
  보이지만 실제로도 작을 수 있다''는 감각(Week 6.1 스케일
  의존성). 교훈은 8.1절 GBDT 학습 곡선과 동일: \textbf{train
  곡선은 ``배운 것을'', val 곡선만 ``재는 것을'' 한다.}}
\end{frame}
